% Pursuing Stacks (cotes 134-2 to 134-8): archive cote 134-6, site page 58 (PDF page 59), typescript page 376.
% typescript. Draft compared: ps5.tex:4267-4339.
% AI-transcribed by Claude Opus 5.5 (claude-opus-5-5, adaptive thinking, effort high), Message Batches API, batch ps3 (msgbatch_01XdFmuVVBAcgc7EFxdtCqsj), 2026-11-12; reviewed against both supplied scan halves; see CHECK_DECISIONS.tsv.
\scanpage{134-6/58}
\mspage{376}

Indeed, by general principles of homological algebra, \unsure{\over[typed]{D}{F}|\over[typed]{F}{D}}or any small
category B, \deleted{due to} \inserted{from} the fact that Addinf(B) is made up with projective
objects of $B^{\hand{\wedge}}_{ab}$ and that any object in $B^{\hand{\wedge}}_{ab}$ is isomorphic to a quotient of
an object in this subcategory, it follows that
\[ (8) \qquad D^{\hand{-}}(B^{\overset{\text{\tikz[baseline=0pt]{\draw[line width=1.1pt] (0,0)--(.18,.12);\draw[line width=1.8pt] (.06,.02)--(.16,.025);\draw[line width=.3pt] (.13,.14) .. controls (.22,.20) and (.25,.13) .. (.17,.10);}}}{\hand{\wedge}}}_{ab}) \hand{\approxeq} W_{B}^{\hand{-}1}\mathrm{Comp}^{\hand{-}}(\mathrm{Addinf}(B)) \ , \]
i.e. the left derived category $D^{\hand{-}}(B^{\hand{\wedge}}_{ab})$ is equivalent with the category
obtained by localizing, with respect to the set $W_{B}$ of \xout{quasi-isomorphisms} \inserted[typed]{homotopy equivalen}%
\inserted[typed]{-ces,}
the category \xout{of} Comp$^{\hand{-}}$(...) of \xout{complexes} differential complexes in the
additive category Addinf(B), with degrees bounded from \deleted{below} \inserted{above} (the diffe-%
rential operator being of degree +1 , according to my preference for
cohomology notation, sorry \hand{!}). An object of $D^{\hand{-}}(B^{\hand{\wedge}}_{ab})$ may thus be viewed
as being essentially the same as \xout{an o} a differential complex in Addinf(B),
with degrees bounded from \deleted{below} \inserted{above}, and given ``up to homotopism'' . The
s\over[typed]{a}{i}m\over[typed]{e}{i}lar description holds for $D^{\hand{-}}(B'^{\hand{\wedge}}_{ab})$ , and in terms of these descrip-%
tions, the ``integration functor'' (in the abelian context) (7) can be
described by
\margin{\tikz[baseline=(lc.base)]{\node[inner sep=3pt] (lc) {\ensuremath{L^{\prime}_{\bullet}}};\draw[color=blue] (lc.center) ellipse (1.8em and 1.6em);}}
\[ (9) \qquad Lf^{ab}_{\hand{!}}(L^{\hand{\prime}}_{\hand{\bullet}}) = \mathrm{Addinf}(L^{\hand{\prime}}_{\hand{\bullet}}) \quad , \]
i.e. by applying componentwise the tautological extension Addinf(f) of f
to the differential complexes in Addinf(B'). This very concrete descripti%
-on applies notably to the complex
\[ (10) \qquad L_{\bullet}^{B'/B} \text{ or } L_{\bullet}^{f} \ \overset{\mathrm{def}}{=}\ f^{ab}_{\hand{!}}(L_{\bullet}^{B'}) \]
introduced yesterday, whenever a (quasi special) integrator $L_{\bullet}^{B'}$ for B'
has been chosen. Applying this to the case of the map
\[ p_{B} : B \hand{\longrightarrow} \hand{\Delta}_{o} \ , \]
we get \xout{an e} (for a given integrat or $L_{\bullet}^{B}$ for B) an explicit description of
the abelianization of the homotopy type of B \xout{i\illegible{3}of} in terms of the
chain complex $p^{ab}_{B\hand{!}}(L_{\bullet}^{B})$ in (Ab) , with the n.th component given by
\[ (11) \qquad \deleted{\xout{\unsure{p^{ab}_{B\hand{!}}}\unsure{(L}_{\illegible{1}}^{\unsure{B}}\illegible{1}}} \ (L_{\bullet}^{B/pt})_{n} \ = \ \underline{Z}^{(I_{n})} \quad , \]
where $I_{n}$ is the set of indices used for describing $L_{n}^{B}$ as the direct
sum of objects of the type $\underline{Z}^{(b_{i})}$ .

Returning to the case of a general \xout{morphism} map f: B'$\hand{\rightarrow}$B , maybe I
should still write down the formula generalizing (2) or (10) of yester-%
day's notes (pages 367 and 369), relating $L_{\bullet}^{B'/B}$ to the cohomology
properties of the map f i.e. to cointegration relative to f. The formula

\margin{58}

% Visual collation: both supplied image halves inspected; see CHECK_DECISIONS_134_7_8.tsv.
